\documentclass{article}
\usepackage[T1]{fontenc}
\usepackage{lmodern}
\usepackage{graphicx}
\usepackage{amsmath,amssymb}
\usepackage[a4paper,margin=2cm]{geometry}
\usepackage[absolute,overlay]{textpos}
\setlength{\TPHorizModule}{1mm}
\setlength{\TPVertModule}{1mm}
\usepackage[indonesian]{babel}
\usepackage{hyperref}
\setlength{\emergencystretch}{2em}
% unit: Halaman 26

\begin{document}

% === Halaman 26 (python/native) ===

4. Forgery Attack

• Tujuan: membuat pesan palsu yang diterima sebagai pesan sah oleh sistem.

• Serangan ini sangat penting pada sistem yang menggunakan MAC (message 

authentication code) atau digital signature (akan dijelaskan pada meteri kuliah 

selanjutnya).

• Misalkan Alice mengirim Pesan + MAC kepada Bob.

• Bob memeriksa MAC untuk memastikan pesan benar-benar berasal dari pihak 

yang memiliki kunci.

• Penyerang, Eve, ingin membuat:

        Pesan palsu + MAC yang valid

   sehingga Bob menganggap pesan tersebut autentik.

26

\end{document}
